\section{Method}
\label{sec:method}

The system is a six-stage pipeline. \textbf{Ingestion} registers a source PDF
and its provenance. \textbf{Extraction} recovers text blocks with their page
numbers and page coordinates, using a hybrid of OCR and the PDF's own text
layer, since either alone is unreliable on the multi-column and tabular layouts
these documents favor. \textbf{Detection} is where structure is recovered,
emitting a flat list of typed elements, each carrying a level, a code, a title,
an optional description, an age band, and the source text it was read from.
\textbf{Parsing} resolves that flat list into the canonical schema of
\S\ref{sec:schema}, attaching each indicator to its full ancestry and composing
its identifier. \textbf{Validation} checks each record against the schema and
against structural invariants on the identifier, and \textbf{persistence}
writes the surviving records to a relational store. The two stages that recover
structure are the subject of this section.

\subsection{Two-Stage Design Rationale}
\label{sec:method-two-stage}

Recovering a standards hierarchy involves two decisions that are easy to
conflate and better kept apart: \emph{what is on the page}, and \emph{what it
is a child of}. The first is local (a heading, a numbered item, a paragraph of
prose) and can be judged from the text in front of the reader. The second is
global: an element's parent may be several pages back, may be printed only once
at the top of a section that runs for multiple pages, or may not be printed at
all and must be recovered from the codes of its own descendants. Detection is
accordingly forbidden from reconstructing an absent parent out of the position
identifiers inside a child's code, and emits the children alone
(Appendix~\ref{rule:emit-what-you-see}); recovering that parent is parsing's
work, not detection's.

I therefore split the work. \textbf{Detection} reads the document in
overlapping chunks and emits a flat list of typed elements, each labeled with
its hierarchy level but not linked to a parent. \textbf{Parsing} takes that flat
list, groups it by domain, and resolves each indicator's full ancestry and
identifier. The split buys three things. It lets each stage see the context it
actually needs, so parsing reasons over an entire domain's elements at once
rather than over one chunk's worth of text. It lets the two stages be graded
independently against separately annotated goldens, which is what makes a
detection failure distinguishable from a resolution failure. And it lets the two
stages run on different models, since they are not equally hard
(\S\ref{sec:method-model-assignment}).

\subsection{Detection}
\label{sec:method-detection}

Detection is two passes over the extracted text blocks.

\paragraph{Pass 1: depth-map inference.} Before any element is classified, a
single cheap model call infers the document's \emph{depth map}: how many nesting
levels this document uses, and what an element at each level looks like in it.
This is the paper's central move. Rather than asking the model to recognize
``strand'' or ``benchmark'' as level names, vocabulary that no two states
share, I ask it once, per document, to characterize that document's own
nesting scheme, and then treat the answer as authority for every subsequent
classification. Appendix~\ref{sec:prompt-depthmap} gives that prompt in full;
it is short, and it is the only one reproduced whole.

The pass runs on a sample of blocks rather than the whole document, under a
6{,}000-token budget, and how that sample is drawn turns out to matter a great
deal. Because every element's level is assigned against the map, \emph{a level
missing from the sample is a level missing from the entire run}: if the sample
contains no evidence that a document nests four deep, the map reports three, and
every element below the missing level is shifted up one.

The obvious sampler, keeping every $n$th block, is the wrong one here. It
keeps each line with the same probability $1/n$ however much structural evidence
that line carries, and in a standards document evidence and frequency run
opposite to one another: a section heading appears once per section, while the
body prose beneath it is repeated on every page. Uniform sampling therefore
drops precisely the headings that establish the top of the hierarchy and keeps
the prose, which establishes nothing, on volume alone.

I stratify by layout instead. Extraction already records each block's normalized
left edge, and that edge is already part of what the model is shown; since
indentation tracks nesting, a rare horizontal position \emph{is} a rare depth.
Blocks are therefore bucketed by left edge, and each bucket is guaranteed a
floor of twelve blocks spread across the document. Buckets are filled
rarest-first, so when the budget binds, what it truncates is the bulk of the
document rather than the structurally distinctive levels. The guarantee is cheap
for the same reason it is needed: those levels are rare, and in Kentucky's
complete standards document the heading buckets together hold only a small
fraction of all blocks.\footnote{An observation from the project's engineering
notes rather than a recorded experiment.} Blocks whose page geometry is unusable
have no left edge to group on, and fall back to a coarse text-shape key reading
word count and final character only. Either way the rule reads layout position
and counts, never any document's words.

\paragraph{Pass 2: classification against the map.} The document is then chunked
into roughly 2{,}000-token windows with 500 tokens of overlap, and each chunk is
classified with the depth map supplied as context. The overlap exists because
elements straddle chunk boundaries; because chunks overlap, a long description
split across a boundary can be reconstructed by splicing the shared span, which
recovers passages that neither chunk holds whole.

The prompt states general document-structure principles rather than per-state
rules. Its central instruction is positional: classify an element by where it
sits in the depth map, not by the word its label happens to use
(Appendix~\ref{rule:positional}). Supporting
rules are stated at the same level of generality: a heading of the form
\texttt{<Label> <id>: <Title>} has the label-and-id as its code and the text
after the colon as its title, for \emph{any} structural label word rather than a
fixed list; a code, when the document prints one, is authoritative and may
appear inline, in a caption beside the heading, or as the shared leading prefix
of its descendants' codes; and when the document prints none, a code is
derived from the title by a stated abbreviation procedure capped at five
characters. All three are one rule, reproduced at
Appendix~\ref{rule:code-resolution}.

\paragraph{Where determinism is enforced in code.} The last of those rules,
deriving a code from the element's title when the document prints none, is
different in kind from the others. The rest ask the model for a judgement about
a document; this one asks it to execute a deterministic string algorithm, and a
sampling model executes such an algorithm inconsistently. In an engineering
measurement on Kentucky (three runs at temperature zero against one frozen
extraction, recorded in the project's design notes rather than as a paper
artifact) roughly a quarter of the elements received a different code on at
least one run, while level, source page, and age band never varied. Because the
code becomes part of a database primary key, a rate near zero is not good
enough. The abbreviation is therefore also executed in Python, but only where
the document supplied no code of its own, decided by two independent guards: a
\emph{shape} guard, since the abbreviation branch emits uppercase letters and
nothing else, so a code carrying a digit, separator, or lowercase letter is
positional and is never touched; and a \emph{grounding} guard, since a code read
off the page appears in the element's own source text while an invented
abbreviation does not. Both are required: a model that transcribes a lettered
item without its list prefix defeats grounding alone.

This is a pairing, not a substitution. The rule stays in the prompt verbatim;
the code makes its floor zero. The same pattern recurs throughout the system,
and \S\ref{sec:method-llm-first} states the discipline that governs it.

\subsection{Parsing}
\label{sec:method-parsing}

Parsing consumes the flat element list, batched by domain so that an entire
domain's elements are resolved together, and produces normalized standards.

The substantive work is composing each indicator's fully-qualified code from its
chain of ancestors, and the difficulty is that printed code namespaces are not
uniform. Three shapes recur. A namespace may \emph{skip} a level, as Nevada's
does: its sub-strand code does not extend its strand code, and requiring it to
would corrupt every one of the 24 annotated Nevada standards, all of which break
the chain at exactly that level. A sub-level identifier may be
\emph{dotted}, restating its parent's position inside its own id, as Kentucky's
\texttt{Benchmark 1.1} does, whose leading \texttt{1} \emph{is} Standard~1; so
composing it onto the full parent code counts that parent twice. And a heading
may print a label-and-id form that must be converted into the domain-qualified
form before it can be composed at all (Appendix~\ref{rule:label-identifier}).

Each is handled by a rule keyed on structure rather than vocabulary. Peeling an
ancestor's code from a descendant's stops where the namespace stops: if a level's
anchored code would equal the code of the level directly below it, the peel has
run past the end of the namespace and that level keeps its own heading's
identifier (Appendix~\ref{rule:namespace-skip}). A code that already carries its
own qualification is used as it stands and is never qualified a second time, at
any level (Appendix~\ref{rule:already-qualified}). A dotted child whose added segment repeats its parent's last segment
is collapsed, but only when the added tail is itself dotted, since a single added
segment is an ordinary child index. A repair that cannot demonstrably improve the
row's own ancestry is not made, and a repair that would introduce whitespace into
an identifier declines rather than trading one malformation for another. None of
these rules contains a list of label words; whitespace and dot structure are the
only signals they read.

Uniqueness is enforced after the merge rather than during it, because a collision
can span batches. Colliding standards are re-qualified using their own ancestors'
segments, and \emph{every} member of a colliding set is rewritten, including the
first one seen, so the result does not depend on the order in which chunks
arrived. Where two standards share every ancestor and remain indistinguishable, a
numeric suffix is assigned in document order. That is reproducible, but it
carries no information, which is why it is a last resort and is logged as one.

Finally, a validation stage rejects any record whose codes are structurally
impossible before it reaches the database: no code may contain whitespace, the
indicator's code must extend its nearest present ancestor's, and the identifier
must end with the indicator code. The middle condition is deliberately scoped to
the leaf, because, as Nevada shows, an intermediate level may legitimately
break the chain.

\subsection{Model Assignment}
\label{sec:method-model-assignment}

The three model calls in the pipeline are not equally difficult, and I assign
the cheapest model that suffices for each. The depth-map pass produces a short
structural summary from a sampled fraction of the document and runs on Claude
Haiku 4.5. Detection is the hard classification problem and runs on Claude Opus
4.6. Parsing operates on already-typed elements with the ancestry constraints
supplied explicitly, and runs on Claude Sonnet 4.6. Since detection is the only
stage invoked once per chunk of the whole document, this assignment also
concentrates spend where the difficulty is. I state the principle here as a
design commitment; \S\ref{sec:experiments-results} reports what the assignment
costs, and a model-tier ablation is identified as future work rather than
claimed as evidence.

\subsection{Confidence and Human Verification}
\label{sec:method-confidence}

Every detected element carries a confidence score, and \textbf{nothing in the
pipeline thresholds it}. There is no review gate, no minimum score, and no
field marking an element as requiring attention; all elements flow through
detection, parsing, validation, and persistence regardless of score. The prompt
does ask for a band, and Appendix~\ref{rule:confidence-bands} reproduces what it
asks for; nothing downstream reads the answer. I am
explicit about this because it is easy to assume otherwise, and because my own
measurement of the score's distribution argues against building such a gate: the
scores occupy five distinct values in $[0.90, 0.97]$, the band the prompt
describes as indicating uncertainty is used zero times, and the score varies
more by document than by whether an element is correct
(\S\ref{sec:experiments-confidence}).

Human verification is a separate mechanism. The schema carries per-element
verification state at all four levels, and an editorial interface lets a
reviewer confirm or correct each element against the source PDF. It is a
workflow over the stored corpus, not a filter applied during extraction, and it
is blanket rather than score-triggered.

Whether every standard must be verified before the corpus is useful is
therefore a cost-and-trust decision rather than a precondition. The durable win
is that fifty independently-formatted documents become one queryable corpus,
and that value is available before every leaf has been signed off, provided the
verification state travels with the data so a consumer can decide what it
requires: uses carrying regulatory or assessment weight will reasonably demand
verified records, exploratory ones reasonably will not.

\subsection{Deployment}
\label{sec:method-deployment}

The method is deployed as a serverless pipeline whose stages communicate by
writing JSON to S3 rather than by passing payloads directly, so each stage's
input and output can be inspected in isolation when a run is debugged.
Detection and parsing are batched through a prepare $\rightarrow$ \texttt{Map}
$\rightarrow$ merge pattern with at most three concurrent workers, so that a
long document does not exhaust a function timeout; detection batches by
text-block chunk and parsing by domain, which keeps related elements in the
same call. Appendix~\ref{sec:appendix-architecture} gives the deployed
architecture and the interface through which the extracted corpus is reviewed.

\textbf{The evaluation suites exercise the direct path rather than this batched
one, and the two are known to diverge}
(\S\ref{sec:discussion-limitations}). Every subset-tier quality number in this
paper is therefore a measurement of the direct path;
\S\ref{sec:experiments-scale} and \S\ref{sec:experiments-scale-quality} are the
evidence that the batched path behaves as intended.

\subsection{LLM-First Design Discipline}
\label{sec:method-llm-first}

The system holds an explicit rule about where structural decisions may live:
\emph{in the prompt, as a general document principle; not in Python, as a
per-document branch}.

This was adopted in response to a measured failure. An earlier revision achieved
strong scores on the states it had been developed against by accumulating
targeted logic for each: regular expressions for particular label words, a
special case for one state's collision behavior, an abbreviation helper tuned to
another's conventions. The scores were real and the generalization was not. A
migration removed that logic and re-expressed each rule as a general principle in
the prompt: instead of a pattern matching \texttt{Strand N:} or
\texttt{Concept N:}, a rule about what a \texttt{<Label> <id>: <Title>} heading
means for any label word; instead of a per-state routine stripping a trailing
structural noun from domain titles, a rule that such a noun is not part of the
name.

Deterministic code is retained, but only where it earns its place by reading the
\emph{shape} of an output rather than any document's vocabulary: JSON extraction
and schema validation; identifier derivation from already-clean fields;
cross-chunk reconciliation of codes that disagree between chunks; the
canonicalization of absence to a single spelling; the layout-stratified sampler;
the code repairs of \S\ref{sec:method-parsing}. Each is paired with a prompt rule
rather than substituted for one, and each exists because a model sampling at
temperature zero emits both spellings of some decision intermittently. A
prompt rule lowers that rate but cannot drive it to zero, and a primary key needs
zero.

The discipline is enforced, not merely stated. A new per-state pattern in the
detector or parser is treated as a regression even when it raises a score on the
development states, and every change is validated against held-out states before
it is accepted. The rule-based comparison of \S\ref{sec:experiments-baseline} is
built deliberately as a throwaway outside the pipeline for the same reason: it is
the concrete form of the approach this discipline forbids, kept so the difference
can be measured rather than asserted.

\paragraph{What the prompts do carry.}
\label{sec:method-discipline}
The discipline above is a claim about \emph{code paths}, and it should not be
read as a claim that the prompts are free of the corpus. They are not. Both
prompts state their rules generally, for any structural label word and at any
nesting depth, but each rule is illustrated with worked examples, and those
examples were written while looking at real documents. Checking every annotated
golden title of at least twelve characters against the verbatim prompt source
finds fifteen that are distinctive to a single state's golden, plus one domain
name that several states share and that therefore carries no
provenance.\footnote{\raggedright Regenerated by
\texttt{paper/\allowbreak analysis/\allowbreak prompt\_\allowbreak
example\_\allowbreak provenance.py}; record at
\texttt{paper/\allowbreak results/\allowbreak task12\_20260905/}.}
Ten of the distinctive strings come
from the four development states, which is ordinary few-shot practice.

Five do not, and I state that plainly because it qualifies how the held-out
states should be read. Four Kentucky strings and one Nevada string appear in the
prompts, and two of the Kentucky ones are long indicator sentences
(``Engages in an activity for a sustained period of time'' $\rightarrow$
\texttt{EASPT}) rather than generic vocabulary. They illustrate the
abbreviation procedure of \S\ref{sec:method-detection}, which is a deterministic
string algorithm over a title and reads no document's meaning; nothing about a
level's classification, a code's namespace or a document's depth was written
from those states. But the honest description of Kentucky and Nevada is
therefore \emph{not annotated during development, and never fitted to} rather
than \emph{never seen}, and the held-out results in
\S\ref{sec:experiments-headline} should be read with that distinction in view.
